\section{Experimental Setup}
\label{sec:setups}

\subsection{Data}
The main source to generate research ideas is the scientific literature $\mathcal{L}$, which we obtain from the Semantic Scholar Academic Graph API\footnote{\href{https://www.semanticscholar.org/product/api}{https://www.semanticscholar.org/product/api}}. From this, we select papers appearing after May 01, 2023, because LLMs that we use in our experiments are trained on data from the open web available before this point. This follows the procedure of existing literature-based hypothesis generation work~\cite{zeroshot/hypothesis}. Then, we select high-impact papers (that have more than 20 citations) as core papers, mirroring human researchers' tendency to leverage influential work, to ensure the high quality of generated ideas. The resulting data is still very large; thus, we further sample a subset of 300 papers as core papers to obtain a reasonably sized benchmark dataset. The average number of reference papers for each core paper is 87; the abstract of each paper has 2.17 entities on average. The distribution of disciplines for all papers is provided in Figure~\ref{fig:discipline}.

\subsection{Baselines and Our Model}
As we target the novel task of research idea generation involving the generation of problems, methods, and experimental designs, there are no baselines for direct comparison. Thus, we mainly compare our full ResearchAgent model against its ablated variants, outlined as follows:
\begin{enumerate}[itemsep=1.0mm, parsep=1pt, leftmargin=*]
    \item {\bf Naive ResearchAgent} -- which uses only a core paper to generate research ideas. 
    \item {\bf ResearchAgent w/o Entity Retrieval} -- which uses the core paper and its relevant references without considering entities.
    \item {\bf ResearchAgent} -- which is our full model that uses the relevant references and entities along with the core paper, to augment LLMs.
\end{enumerate}
In addition to this set of core baselines, we also compare our approach against existing hypothesis generation work from prior literature in Table~\ref{tab:hypothesis}.

\subsection{Evaluation Setup}
Given our formulation of idea generation~(Sec~\ref{sec:research-idea-gen}), there are no ground-truth answers to measure the quality of the generated ideas. Yet, exhaustively listing pairs of core papers and reference research ideas is suboptimal, since there may exist a large number of valid research ideas for each core paper, and this process requires much time, effort and expertise on the part of human researchers. Thus, we use a combination of model-based automatic evaluation and manual human evaluation to validate different models on our experimental benchmark.

\begin{figure*}[t!]
    \centering
    \begin{subfigure}[t]{0.99\linewidth}
        \centering
        \includegraphics[width=0.75\linewidth]{Figures/radar_human.pdf}
        \caption{Human Evaluation}
    \end{subfigure}
    
    \vspace{0.1in}
    
    \begin{subfigure}[t]{0.99\linewidth}
        \centering
        \includegraphics[width=0.75\linewidth]{Figures/radar_model.pdf}
        \caption{Model-based Evaluation}
    \end{subfigure}
    
    \vspace{-0.075in}
    \caption{Main results on our research idea generation task with human- (top) and model-based (bottom) evaluations, where we report the score of each idea (problem, method, or experiment design) based on its own five criteria and their average score.}
    \label{fig:main}
    \vspace{-0.05in}
\end{figure*}

\begin{figure}[t!]
    \centering
    \includegraphics[width=0.975\columnwidth]{Figures/pairwise_fig.pdf}
    \vspace{-0.075in}
    \caption{Results of pairwise comparisons between ideas from two of any different approaches, where we report the win ratio.}
    \label{fig:pairwise}
    \vspace{-0.1in}
\end{figure}

\paragraph{Model-based Evaluation}
Following the recent trends in using LLMs to judge the quality of output texts (especially in the setting of reference-free evaluations)~\cite{LLM/Judge/1, LLM/Judge/2, G-Eval}, we use GPT-4 to judge the quality of research ideas. Note that each of the problem, method, and experiment design is evaluated with five different criteria (See labels of Figure~\ref{fig:main} for criteria and see Table~\ref{tab:criteria} for their detailed descriptions). We ask the LLM-based evaluation model to either rate the generated idea on a 5-point Likert scale for each criterion or perform pairwise comparisons between two ideas from different models. We provide the prompts for evaluations in Appendix~\ref{appendix:setups}.  

\paragraph{Human Evaluation}
Similar to model-based evaluations, we perform human evaluations that involve assigning a score for each criterion and conducting pairwise comparisons between two ideas. As the generated ideas are knowledge-intensive, we carefully select annotators who are well-versed in the field and provide them with ideas that are highly relevant to their field of expertise\footnote{We also experiment with human evaluation using non-domain-experts, but this proves to be suboptimal therefore, we focus on experts for reliable judgments of generated ideas.}. Specifically, we choose ten expert researchers who have authored at least three papers and ask them to judge only the ideas that are generated based on their own papers.

\subsection{Implementation Details}
We mainly use the GPT-4~\cite{GPT-4} release from Nov 06, 2023, as the basis for all models, which is, notably, reported to be trained with data up to Apr 2023 (meanwhile, the papers used for idea generation appear after May 2023). To extract entities and build the entity-centric knowledge store, we use the off-the-shelf BLINK entity linker~\cite{blink}, with papers from May 01, 2023, to Dec 31, 2023 (available from Semantic Scholar API) along with their references, which number 50,091 in total. We provide detailed prompts used to elicit responses for research idea generation in Appendix~\ref{appendix:prompts:generation}.


